%-----------------------------------------------------------------------------
% CONCLUSION
%-----------------------------------------------------------------------------
\section{Conclusions}
\label{sec:conclusions}
\color{black}

This thesis investigated the accuracy-energy trade-off of a multi-agent
LLM system for sales-data analysis. The central result is that model choice
alone is not sufficient to explain performance: the most important effects
emerge from the interaction between model family, agent step, prompt
difficulty, and inference-time parameters.

The baseline comparison shows that \texttt{mistral-small3.2:24b} is the
most efficient untuned model. It achieves the best aggregate strict quality
at the lowest energy cost, especially on easy and medium prompts. However,
the prompt-difficulty analysis reveals a limitation that is hidden by the
aggregate score: Mistral collapses on the hardest SQL aggregation prompts,
where \texttt{gemma4:26b} performs best at baseline and \texttt{gemma4:e4b}
can be repaired through targeted lookup interventions. This confirms the
importance of reporting results by prompt difficulty rather than only as a
single mean score.

The per-agent sensitivity experiments show that hyperparameter tuning must
be step-specific. Gemma~26B is highly tunable but fragile: its visualization
agent is weak at baseline and responds strongly to both static sampling
parameters and larger token caps. In particular, raising the Gemma~26B
visualization cap from 7000 to 10000 tokens improves \texttt{vis\_score}
with negligible energy cost, and the curve does not yet show a clear
plateau. Gemma~E4B behaves differently: adding too much token headroom can
reduce quality, especially in lookup and visualization, so it requires
bounded caps rather than simply larger ones. Mistral is almost insensitive
to token caps; moderate budgets are enough for stable outputs.

The token-budget experiment therefore supports a practical rule: increasing
\texttt{max\_tokens} is cheap when the model does not actually consume the
extra budget, so generous caps are useful for model-step pairs that benefit
from additional headroom. They should not be treated as universally better,
because some thinking-model agents become more verbose and less reliable
when given too much room.

The compute-expansion experiment shows that CoT refinement is useful only
when the feedback signal matches the failure mode. Lookup CoT is effective
because SQL errors are executable and can be fed back to the model. For
Gemma~E4B, lookup CoT is the strongest single intervention in the thesis;
for Mistral it also helps, but the gain is smaller. Visualization CoT is
harmful for both models. The generated chart often fails semantically rather
than syntactically, so code-execution feedback does not reveal whether the
plot matches the user's requested grouping, split, or chart semantics.
Additional CoT iterations therefore tend to over-edit the output and degrade
\texttt{vis\_score}. The corrected early-stop diagnostics reinforce this
point: requested depths of 3 or 4 usually execute only around two iterations,
so high \texttt{cot\_n} should be interpreted as an upper bound rather than
as the actual compute spent. Depth~2 is the safest default, while depth~3 can
be reserved for hard lookup prompts when maximum accuracy is more important
than energy.

The final Pareto confirmation combines the best insights from the earlier
tests. Under GPT-5.4 judging, the final non-dominated frontier is dominated
by low-energy Mistral candidates and one high-quality Gemma~E4B endpoint.
The best final Gemma~E4B configuration combines lookup CoT, a lower lookup
temperature, and efficient downstream token caps; it reaches the highest
strict prompt quality with full completion. Mistral remains the natural
choice for low-energy execution, especially on easy-to-medium prompts.
Gemma~26B improves with tuning, but its energy cost keeps it off the final
full-benchmark Pareto frontier.

Overall, the thesis supports three practical recommendations for LLM data
agents. First, tune parameters per agent step rather than globally. Second,
allocate compute according to the failure mode: token headroom for steps
that are genuinely truncated, CoT for SQL lookup where feedback is
verifiable, and conservative generation for visualization when the main
error is semantic mismatch. Third, report accuracy and energy by prompt
difficulty, because the best model for simple tasks is not necessarily the
best model for hard analytical questions.
